% ==============================================================================
% Fichier : front/02_mode_emploi.tex : SAFI (SIS589)
% ==============================================================================

\chapter*{Mode d'emploi}
\addcontentsline{toc}{chapter}{Mode d'emploi}
\markboth{Mode d'emploi}{Mode d'emploi}

\section*{À qui s'adresse ce cours}

Aux étudiants du Master Professionnel en Sécurité des Systèmes Informatiques
et des Données, et aux professionnels qui exercent ou se destinent aux
métiers d'\textbf{analyste SOC}, d'\textbf{analyste forensique}, de
\textbf{responsable de la réponse à incident} et de \textbf{RSSI}. Des
notions de réseaux (TCP/IP), d'administration Windows et Linux et de droit
numérique sont utiles.

\section*{Comment lire ce livre}

Le livre comprend un prologue et six parties. Chaque chapitre suit le même
plan~: une épigraphe, le niveau et les prérequis, les objectifs, une question
directrice, cinq à neuf sections, un encadré « à retenir », dix exercices
gradués et une transition. Les niveaux des exercices sont signalés par des
étoiles~: \niveaubase{} fondamental, \niveaumoyen{} intermédiaire,
\niveauexpert{} expert.

\section*{Les encadrés}

\begin{itemize}
    \item \textbf{Scénario opérationnel}~: un épisode des opérations BÉNOUÉ,
          FARO ou MAYO, avec \emph{constatation} (faits) et \emph{analyse}
          (interprétation, niveau de confiance).
    \item \textbf{Artefact}~: la fiche d'une pièce d'enquête (nature, source,
          pertinence, fiabilité, opposabilité).
    \item \textbf{Chaîne de custody}~: une procédure de préservation de la
          preuve.
    \item \textbf{Analyse CRO}~: l'arbitrage Confidentialité, Fiabilité,
          Opposabilité.
    \item \textbf{Application métier}~: ce que le chapitre change pour un
          métier.
    \item \textbf{Cadre juridique}~: le texte de droit applicable.
\end{itemize}

\section*{Constatation et analyse}

Le cours distingue systématiquement ce que l'on \textbf{constate} (un fait
vérifiable) de ce que l'on en \textbf{déduit} (une interprétation, assortie
d'un niveau de confiance). C'est la discipline première de la supervision et
du forensique.

\section*{NDAP SUKULU et évaluation}

L'environnement compagnon NDAP SUKULU (page suivante) propose, pour chaque
chapitre, une fiche d'auto-évaluation et, pour les chapitres pratiques, un
lab sur pièces. Fiches et labs se font hors des 30~heures de cours. Le code
du cours est \textbf{\codeUE}.
